%% ai_engineer_resume.tex
%% -----------------------------------------------------------------------------
%% AI Engineer resume -- CONTENT ONLY.
%%
%% All styling lives in resume.cls. Edit the text below to update the resume;
%% do NOT change fonts/spacing/colours here (do that in resume.cls) so the
%% template stays consistent. Compile with:  pdflatex ai_engineer_resume.tex
%% -----------------------------------------------------------------------------
\documentclass{resume}

\begin{document}

% ---- Header ----------------------------------------------------------------
\name{Chinmay Harjai}
\contact{%
  \phone{+91 7878709846} \contactsep
  \email{chinmayharjai@gmail.com} \contactsep
  \github{chinmayharjai}{GitHub} \contactsep
  \linkedin{https://www.linkedin.com/in/chinmayharjai}{LinkedIn}%
}

% ---- Education -------------------------------------------------------------
\section{Education}
\begin{itemize}
  \education{BTech, Electrical Engineering}{National Institute of Technology, Jaipur}{2025}
\end{itemize}

% ---- Experience ------------------------------------------------------------
\section{Experience}
\entry{ICICI Bank \textnormal{— Data Engineer, Mumbai}}{July 2025 -- Present}
\begin{itemize}
  \item Worked with IBM InfoSphere DataStage and Designer to schedule and monitor automated ETL jobs, ensuring reliable batch processing and timely data availability across enterprise systems.
  \item Analyzed and resolved UCIC (Unique Customer Identification Code) mismatches by investigating MDM deduplication rules and API matching logic, improving customer identity accuracy across downstream applications.
  \item Built and enhanced Azure ETL pipelines using Azure Data Factory and Azure Databricks to ingest and transform source-system data, delivering analytics-ready datasets to business and risk teams.
  \item Implemented PII and PCI data redaction using regular expressions to identify sensitive fields across on-premise Oracle and MySQL databases, strengthening data security and compliance controls.
\end{itemize}

% ---- Projects --------------------------------------------------------------
\section{Projects}

\project{Multi-Agent Content Intelligence \& Publishing System}
\begin{itemize}
  \item Designed a multi-agent system using LangGraph where specialised agents (research, drafting, quality-check) coordinate via shared state to autonomously generate and publish LinkedIn content.
  \item Exposed agent orchestration logic through a FastAPI microservice, enabling async job submission, streaming LLM responses, and webhook-based triggering from the n8n workflow layer.
  \item Containerised the full stack with Docker on GCP, managing session persistence, secret injection, and retry logic for production-grade reliability.
\end{itemize}

\project{Conversational AI Assistant with Context Memory (WhatsApp)}
\begin{itemize}
  \item Built a production conversational AI assistant on WhatsApp using WAHA + LangChain, implementing session-based context management and a short-term memory buffer to maintain conversation continuity.
  \item Designed a vector-backed retrieval layer (ChromaDB) for FAQ grounding, reducing hallucinated responses on domain-specific queries.
  \item Deployed an async event-handling backend with FastAPI and Docker on a GCP Linux VM, integrating a rule-based AI classifier to filter unsafe or off-topic responses before delivery.
\end{itemize}

\project{Agentic RAG Pipeline for Quantitative Alpha Research}
\begin{itemize}
  \item Built a RAG pipeline using LangGraph to orchestrate an autonomous research agent that ingests unstructured financial documents, retrieves relevant context via semantic search, and generates structured alpha signals using LLMs.
  \item Built a vector retrieval layer using ChromaDB with Gemini embeddings and cosine similarity search, grounding LLM outputs to reduce hallucination on domain-specific financial data.
  \item Implemented async Python workflows with queue-based parallel execution and SQL logging for experiment tracking; used the WorldQuant Brain API to programmatically submit alpha formulas and retrieve backtesting results.
\end{itemize}

% ---- Technical Skills ------------------------------------------------------
\section{Technical Skills}
\begin{itemize}
  \skill{Programming Languages}{Python, SQL}
  \skill{AI \& ML}{PyTorch, TensorFlow, LSTM, NLP, LangChain, LangGraph, Hugging Face, ChromaDB}
  \skill{Database \& Analytics Tools}{ETL (IBM DataStage), MDM, Azure Data Factory, Databricks}
  \skill{Data Visualisation}{Tableau, Power BI}
  \skill{Software}{Visual Studio Code, GitHub Desktop, MS Office (Excel, Word, PowerPoint), GCP}
  \skill{Others}{Dynatrace, Postman, Spark SQL, SQL Developer, n8n, Docker}
\end{itemize}

% ---- Leadership ------------------------------------------------------------
\section{Leadership}
\begin{itemize}
  \item \textbf{Executive, Think India MNIT Club:} Led the planning and execution of large-scale events with 5,000+ attendees, managing and coordinating a team of 20+ students to ensure smooth operations and delivery.
  \item \textbf{Founder, Unicorn Club (2022):} Built and scaled a community of entrepreneurs, organizing and hosting 30+ talks focused on practical business insights, collaboration, and knowledge sharing.
\end{itemize}

\end{document}
